% -*- coding: utf-8 -*-
\chapter{Marco teórico / Estado del arte}
\label{ch:estado-arte}

Este capítulo no es un recuento de la inteligencia artificial, el big data, los data lakes ni la privacidad. Se limita a tres preguntas: por qué un LLM alucina en texto normativo; qué familias de RAG existen y \textbf{cuál se implementa y evalúa} en este TFM; y qué han medido trabajos previos sobre alucinaciones, métricas y dominios regulados. La revisión es crítica: se señala qué se puede transferir al PAC/POSEI y qué afirmaciones absolutas no se adoptan.

\section{Arquitectura Transformers y Grandes Modelos del Lenguaje}
Los grandes modelos de lenguaje (LLMs, por sus siglas en inglés) constituyen una de las principales líneas de avance en el Procesamiento del Lenguaje Natural (PLN) debido a su capacidad para modelar dependencias complejas entre tokens y generar texto coherente en múltiples tareas de comprensión y razonamiento. Su desarrollo contemporáneo se fundamenta en la arquitectura Transformer, introducida por Vaswani et al. (2017), la cual desplazó a los modelos secuenciales previos al proponer un mecanismo de atención capaz de representar relaciones contextuales de forma altamente paralelizables.

\begin{figure}[htbp]
\centering
\includegraphics[width=0.62\textwidth]{figures/transformer-vaswani.png}
\caption{Arquitectura Transformer (codificador--decodificador) con atención multi-cabeza, conexiones residuales y normalización (Vaswani et al., 2017).}
\label{fig:transformer}
\end{figure}

\begin{figure}[htbp]
\centering
\includegraphics[width=0.58\textwidth]{figures/evolucion-llm.png}
\caption{Línea temporal de modelos de lenguaje, desde n-gramas y RNN/LSTM hasta Transformer, BERT, GPT y los LLM contemporáneos.}
\label{fig:evolucion-llm}
\end{figure}

A diferencia de las Redes Neuronales Recurrentes (RNN) o las redes de Memoria a Corto y Largo Plazo (LSTM), que procesan la información de manera estrictamente secuencial y sufren de degradación de la señal en secuencias largas (el problema del desvanecimiento del gradiente), el Transformer permite ponderar simultáneamente la interacción entre todos los elementos de una secuencia. Esta arquitectura asume un rol fundacional en este trabajo, ya que constituye el motor inferencial generativo del sistema; sin embargo, comprender sus fundamentos matemáticos es indispensable para identificar por qué, estructuralmente, estos modelos son propensos a la alucinación factual en corpus normativos como la Política Agraria Común (PAC).

La Tabla~\ref{tab:nlp-arch} no trata RAG como una red neuronal equivalente a RNN, LSTM o Transformer: compara \emph{enfoques} de PLN. RNN, LSTM y Transformer son arquitecturas de red; RAG es un patrón de acceso a conocimiento externo que se monta \emph{sobre} un Transformer.

\begin{table}[htbp]
\centering
\caption{Comparación de arquitecturas y enfoques NLP.}
\label{tab:nlp-arch}
\begin{tabular}{p{3.2cm}p{5.4cm}p{5.4cm}}
\toprule
\textbf{Enfoque} & \textbf{Ventaja} & \textbf{Limitación} \\
\midrule
RNN & Memoria secuencial de tokens & Desvanecimiento del gradiente \\
LSTM & Dependencias más largas que la RNN clásica & Lentitud y coste de entrenamiento \\
Transformer & Paralelismo y ventana de contexto amplia & Alto coste computacional; alucinación paramétrica \\
RAG (enfoque, no red) & Conocimiento actualizable y citable en inferencia & Depende de la calidad de la recuperación \\
\bottomrule
\end{tabular}
\end{table}

\begin{figure}[htbp]
\centering
\includegraphics[width=0.78\textwidth]{figures/desvanecimiento-gradiente.png}
\caption{Desvanecimiento del gradiente en retropropagación: la magnitud de la señal decae al alejarse del instante actual, limitando a las RNN clásicas.}
\label{fig:gradiente}
\end{figure}

\subsection{Fundamentos matemáticos de la atención (por qué no hay verificación externa)}
El núcleo del Transformer es la autoatención de producto escalar escalado (Vaswani et al., 2017). Dada una secuencia de embeddings $X$, el modelo construye las matrices de consulta, clave y valor
\begin{equation}
Q = XW_Q,\qquad K = XW_K,\qquad V = XW_V,
\end{equation}
y calcula
\begin{equation}
\mathrm{Attention}(Q,K,V)=\mathrm{softmax}\!\left(\frac{QK^{\top}}{\sqrt{d_k}}\right)V.
\end{equation}
Cada token generado es una mezcla contextual del texto precedente. Esa formulación explica la fluidez y, a la vez, la carencia normativa: $V$ proviene de pesos de preentrenamiento, no de un expediente. No hay, en esta ecuación, un mecanismo que contraste la correlación estadística contra una base documental.

En inferencia decoder-only (LLaMA, GPT), la generación es autorregresiva: el token $t$ se muestrea de $P(x_t\mid x_{<t})$. RAG no cambia esa ecuación; cambia el condicionamiento, al concatenar fragmentos recuperados a $x_{<t}$.

\begin{figure}[htbp]
\centering
\includegraphics[width=0.82\textwidth]{figures/matriz-atencion.png}
\caption{Matriz de autoatención sobre la oración «El banco de madera estaba en el parque»: cada consulta pondera las claves del contexto. Ilustra por qué el sentido de un token depende de su entorno y no de un embedding estático.}
\label{fig:atencion}
\end{figure}

\begin{figure}[htbp]
\centering
\includegraphics[width=0.78\textwidth]{figures/hiperplano-decision.png}
\caption{Hiperplano de decisión lineal $W\cdot X+b=0$ en un espacio de dos atributos. Sirve de analogía geométrica a la separación de representaciones, previa a los espacios densos de los embeddings.}
\label{fig:hiperplano}
\end{figure}

\subsection{Limitaciones Estructurales de los LLMs}
La aplicación directa de Modelos de Lenguaje de Gran Escala aislados (standalone) para la consulta y gestión de documentación técnica o regulatoria presenta severas inconsistencias funcionales y metodológicas. Aunque los LLMs exhiben una notable capacidad para comprender el lenguaje natural, resumir textos complejos y generar respuestas gramaticalmente fluidas, su arquitectura subyacente impone barreras insuperables cuando se exige rigor jurídico, actualización documental estricta y trazabilidad en el ámbito agrario (por ejemplo, en la regulación de las ayudas de la Política Agraria Común ---PAC--- o el Programa de Opciones Específicas por la Lejanía y la Insularidad ---POSEI---).

La diferencia fundamental entre un sistema de información documental y un LLM aislado reside en \emph{dónde} está el conocimiento:

\begin{itemize}
  \item \textbf{Memoria paramétrica.} Todo el conocimiento del LLM reside de forma distribuida e implícita en los pesos de atención y de las redes \emph{feed-forward}, fijados en el preentrenamiento. En inferencia el modelo no consulta un expediente: calcula una distribución sobre el vocabulario.
  \item \textbf{Memoria no paramétrica.} El corpus externo (PDF, tablas, fragmentos indexados) no está en los pesos. RAG la usa en tiempo de consulta: recupera evidencia y la inyecta en el contexto.
\end{itemize}
\subsection{Taxonomía del Problema de las Alucinaciones en Dominios Normativos}
Se define la alucinación en los LLMs como la generación de texto que es fácticamente incorrecto, desalineado del contexto o carente de sustento en la evidencia disponible (Huang et al., 2025; Ji et al., 2023). En el procesamiento de normativas agrícolas como la PAC o el POSEI, las alucinaciones se categorizan estructuralmente en dos tipos:

Alucinaciones Intrínsecas: La salida del modelo contradice explícitamente la información proporcionada en el contexto o en el prompt de entrada. Ocurre cuando el mecanismo de atención pondera erróneamente tokens de contexto contiguos, pero semánticamente antagónicos.

Alucinaciones Extrínsecas: El modelo genera afirmaciones que no pueden ser verificadas ni refutadas a partir del contexto provisto (o cuando se le consulta de forma directa sin contexto). El LLM fabrica datos, artículos de ley inexistentes o requisitos de elegibilidad plausiblemente redactados, pero jurídicamente nulos.

En el marco de la regulación agrícola, una alucinación extrínseca ---como inventar una excepción normativa en el plazo de presentación de solicitudes de pago por ecorregímenes--- conduce a la invalidación del sistema, ya que invalida el criterio de seguridad jurídica y puede derivar en sanciones económicas directas para los solicitantes.

\subsection{Restricciones en la Ventana de Contexto y el Fenomeno Lost in the Middle}
Un intento ingenuo de resolver la falta de conocimiento paramétrico en un LLM standalone consiste en inyectar el texto normativo completo (cientos de páginas del Boletín Oficial del Estado o reglamentos comunitarios de la UE) dentro de la ventana de contexto (context window) de un modelo con capacidad extendida (ej. 32k, 128k o más tokens). Sin embargo, esta estrategia colapsa por dos razones estructurales:

Complejidad Computacional: El mecanismo de autoatención  escala de forma cuadrática  con respecto a la longitud de la secuencia , elevando exponencialmente la latencia y la demanda de memoria VRAM.

Degradación del Posicionamiento Semántico (Lost in the Middle): Como demostraron Liu et al. (2023), los modelos basados en Transformer muestran una marcada asimetría en la capacidad de recuperación dentro del contexto de entrada. Los modelos son altamente eficaces identificando e integrando información situada al inicio (primacy effect) o al final (recency effect) del prompt, pero sufren un menoscabo drástico en el rendimiento cuando la información clave se halla inmersa en la zona central del texto inyectado. En normativas de gran volumen, los condicionantes técnicos clave suelen estar sepultados en anexos intermedios, quedando efectivamente ``invisibilizados'' para la atención del modelo.

\begin{figure}[htbp]
\centering
\includegraphics[width=0.82\textwidth]{figures/lost-in-the-middle.png}
\caption{Fenómeno \emph{Lost in the Middle}: la retención de un dato clave cae en el centro del prompt y se recupera en los extremos (primacía y recencia).}
\label{fig:lost-middle}
\end{figure}

\subsection{Volatilidad Normativa y Falta de Determinismo y Trazabilidad}
Finalmente, la gestión de la normativa agraria exige resolver tres factores críticos incompatibles con los LLMs aislados:

Volatilidad y Asincronía Temporal: Las leyes agrícolas están sujetas a modificaciones periódicas (órdenes de modificación de la PAC, resoluciones autonómicas de sequía, ajustes en las fichas del POSEI). Un LLM standalone es estático; su conocimiento termina de forma tajante en su fecha de corte de entrenamiento (knowledge cutoff). La única vía para actualizar la memoria paramétrica de un LLM es mediante el reentrenamiento o el Ajuste Fino (Fine-Tuning), procesos prohibitivos en términos de tiempo, costes computacionales y riesgo de catastrophic forgetting (olvido catastrófico de habilidades previas).

Ausencia de Trazabilidad y Citación Granular: Un LLM aislado opera como una «caja negra» probabilística. Al generar una respuesta sobre un requisito técnico, no es capaz de proporcionar la dirección exacta de la fuente (p. ej., Reglamento (UE) 2021/2115, Artículo 14, apartado 2, página 45). La cita producida por un LLM aislado suele ser una alucinación estilística basada en patrones de citación aprendidos.

Falta de Determinismo: Dado que el proceso de muestreo en la capa de salida (decoding) emplea parámetros como la temperatura () o el muestreo top-, la misma consulta formulada en momentos distintos puede derivar en interpretaciones contradictorias sobre un mismo marco regulatorio.

\section{Arquitectura Retrieval-Augmented Generation (RAG)}
La arquitectura Retrieval-Augmented Generation (RAG), introducida formalmente por Lewis et al. (2020), nace como un paradigma híbrido que desacopla el almacenamiento del conocimiento de la capacidad de razonamiento del modelo de lenguaje. En lugar de confiar exclusivamente en la \textbf{memoria paramétrica} (los pesos del LLM), RAG integra un almacén \textbf{no paramétrico}: un componente de recuperación semántica que extrae evidencia documental en tiempo de inferencia y la inyecta en el contexto del generador.

En el ámbito de la documentación normativa agraria (PAC/POSEI), esta arquitectura no se concibe como un mero sustituto del LLM, sino como un mecanismo de anclaje estricto (Strict Grounding) que condiciona la probabilidad de generación del modelo a la evidencia recuperada, mitigando estructuralmente las alucinaciones extrínsecas y posibilitando la trazabilidad exacta de las respuestas.

\subsection{Descomposición del Pipeline de RAG: Indexación, Recuperación y Generación}
El funcionamiento de un sistema RAG en entornos regulados se divide operativamente en dos fases principales: Indexación Asíncrona (Ingesta) y Pipeline de Inferencia Síncrona (Recuperación y Generación).

\begin{figure}[htbp]
\centering
\includegraphics[width=0.95\textwidth]{figures/rag-pipeline-simple.png}
\caption{Pipeline RAG: la pregunta dispara una consulta de recuperación; los textos recuperados se concatenan al prompt completo y el LLM genera la respuesta anclada.}
\label{fig:rag-pipeline}
\end{figure}

En este trabajo, la arquitectura RAG no se interpreta como una solución automática al problema de las alucinaciones, sino como un mecanismo cuya eficacia debe ser evaluada empíricamente. La literatura reciente sobre RAGAS distingue explícitamente entre métricas de recuperación y métricas de generación, como context precision, context recall, faithfulness y answer relevancy, precisamente porque la utilidad de un sistema RAG no puede inferirse solo a partir de la fluidez de la respuesta. Por ello, el presente estudio adopta una visión cuantitativa en la que la recuperación semántica y la generación se analizan por separado y en conjunto.

Fase de Indexación y Preparación Documental:

Parseo y Extracción de Maquetación (Layout Analysis): Los documentos normativos en formato PDF se procesan mediante herramientas avanzadas (e.g., Docling con soporte OCR) que preservan la estructura jerárquica, tablas y elementos visuales, exportándolos a sintaxis Markdown estructurada.

Fragmentación Semántica (Chunking): El texto unificado se segmenta en bloques discretos . Para evitar la rotura de contexto normativo en medio de un artículo o tabla, se aplican estrategias de fragmentación por encabezados o Summary-Augmented Chunking (Reuter et al., 2025).

Incrustación Vectorial (Embedding): Cada fragmento  se pasa por un modelo codificador (Encoder) para proyectarlo en un espacio vectorial denso de dimensión :

Los vectores resultantes, junto con metadatos de gobernanza (tenant\_id, página, sección, checksum), se almacenan en un índice espacial.

La calidad de esa proyección condiciona todo el retriever denso: si dos redacciones distintas de la misma norma (por ejemplo, ``ayuda al arranque de viñedo'' y ``prima por abandono de viña'') no quedan próximas en el espacio, el generador nunca verá el artículo correcto. Por ello conviene distinguir las familias de embeddings que hoy dominan la recuperación semántica y justificar cuál se instanció en el prototipo.

\textbf{Sentence Transformers.} Sentence Transformers es una familia de modelos basados en BERT que generan embeddings a nivel de oración o párrafo (Reimers y Gurevych, 2019). Frente al BERT original, que produce representaciones de token y no estaba entrenado para comparar frases con similitud coseno, estos modelos se afinan con aprendizaje siamés o tripletas de modo que la distancia entre vectores refleje similitud semántica. Variantes como \texttt{all-MiniLM-L6-v2} (384 dimensiones, muy ligera) o \texttt{paraphrase-multilingual-mpnet-base-v2} (768 dimensiones, mejor cobertura idiomática) se usaron de forma masiva como primer retriever denso en prototipos RAG. Su ventaja es la eficiencia y el ecosistema Hugging Face; su límite, en un corpus normativo en español, es que los checkpoints más citados están optimizados para parafraseo general o para inglés, y degradan cuando la consulta es una paráfrasis agraria y el documento es un artículo del POSEI con identificadores alfanuméricos.

\textbf{BGE (BAAI General Embeddings).} BGE, desarrollada por la Academia de Inteligencia Artificial de Pekín, reentrenó el paradigma de embedding para búsqueda: instrucciones de consulta, normalización L2 y evaluación sistemática en MTEB (Xiao et al., 2023). El modelo de referencia para escenarios no anglófonos no es un hipotético \texttt{bge-large-es} ---que no forma parte del catálogo oficial de BAAI--- sino \texttt{BAAI/bge-m3}: 1024 dimensiones, más de 100 idiomas, ventanas de hasta 8192 tokens y capacidad de devolver a la vez vector denso, pesos dispersos y multi-vector tipo ColBERT (Chen et al., 2024). En este trabajo BGE no se descartó: se reservó para la segunda etapa del pipeline, el Cross-Encoder \texttt{BAAI/bge-reranker-v2-m3}, que reordena los candidatos ya fusionados por RRF. Así se aprovecha la fortaleza de BGE en ranking fino sin duplicar el coste de reindexar todo el corpus con su encoder bi-encoders.

\textbf{E5 (Embeddings from bidirEctional Encoder Representations).} E5, de Microsoft, introduce un entrenamiento débilmente supervisado y, sobre todo, una asimetría explícita entre consulta y documento (Wang et al., 2022). En inferencia se prefija \texttt{query:} o \texttt{passage:}, de modo que el espacio vectorial no trata la pregunta del inspector y el párrafo del boletín como el mismo tipo de objeto. \texttt{intfloat/multilingual-e5-large} (1024 dimensiones) ofrece un compromiso sólido en español y en recuperación asimétrica. El inconveniente operativo para este prototipo es que E5 vive de forma natural en Hugging Face Transformers, no en el runtime Ollama con el que se sirve el generador local; adoptar E5 habría partido la inferencia en dos stacks (Transformers para embedder y Ollama para LLM).

\textbf{Nomic Embed (elección por defecto del prototipo).} El embedder servido en la demostración contenedorizada es \texttt{nomic-embed-text} (\texttt{RAG\_EMBEDDING\_MODEL} en \texttt{.env.example} y Docker Compose): un modelo de sentencia ligero, con vectores de 768 dimensiones, pensado para recuperación y servido por Ollama junto al generador. Esa elección privilegia reproducibilidad y coste de arranque (CPU o GPU modestas) frente a un embedder 8B. La familia Qwen3 Embedding permanece en el catálogo de contraste del laboratorio (\texttt{qwen3-embedding:latest}) para experimentos embedding $\times$ distancia; no es el valor por defecto de la demo. Los vectores Nomic se almacenan en una columna pgvector de 4096 dimensiones con relleno de ceros, restricción del esquema heredado.

\begin{table}[htbp]
\centering
\caption{Justificación de \texttt{nomic-embed-text} (demo) frente a Sentence Transformers, BGE, E5 y Qwen3, y efecto esperado sobre la recuperación semántica.}
\label{tab:embedding-eleccion}
\small
\begin{tabular}{p{3.1cm}p{6.2cm}p{5.5cm}}
\toprule
\textbf{Criterio} & \textbf{Por qué se eligió este modelo} & \textbf{Impacto en la recuperación semántica} \\
\midrule
Idioma del corpus & Español técnico PAC/POSEI. Nomic y Qwen3 cubren varios idiomas; MiniLM y \texttt{bge-large-en} están sesgados al inglés. & Acerca paráfrasis del agricultor o del gestor (``prima de arranque'') al tenor literal del boletín, elevando recall@k del retriever denso. \\
\addlinespace
Runtime del prototipo & Inferencia local unificada en Ollama, junto al LLM. E5 y Sentence Transformers exigirían un stack Hugging Face paralelo. Nomic arranca en CPU; Qwen3-8B exige más VRAM. & Misma API de embeddings y reindexado reproducible sin servicios cloud. \\
\addlinespace
Asimetría consulta/documento & Qwen3 y E5 son \emph{instruction-aware}; Nomic trata query y passage de forma más simétrica. & El ranking denso de la demo es más simple; el laboratorio permite contrastar Qwen3 cuando hay GPU. \\
\addlinespace
Dimensionalidad & Nomic: 768 dimensiones (columna pgvector de 4096 con padding). Qwen3 llega a 4096; MiniLM usa 384; E5/BGE-M3, 1024. & Espacio suficiente para la demo; el padding es una limitación de esquema, no una ganancia semántica. \\
\addlinespace
Ventana de contexto & Los chunks del prototipo son de 1400 caracteres con solape 120. MiniLM clásico se queda en 512 tokens. & El embedding del fragmento captura el artículo o la ficha, no un recorte a 512 tokens que parte tablas LMR. \\
\addlinespace
Papel de BGE en el sistema & BGE no se usa como embedder por defecto; el Cross-Encoder \texttt{bge-reranker-v2-m3} está cableado pero \texttt{RAG\_USE\_RERANKER=false} en la demo. & La recuperación semántica de la demostración es denso Nomic + BM25 + RRF, sin rerank. \\
\addlinespace
Complemento léxico & El denso no sustituye a BM25: los códigos (``Orden APA/102/2024'') siguen yendo al índice léxico. & Menos document-level mismatch: sinonimia (denso) + identificador exacto (BM25) + RRF. \\
\bottomrule
\end{tabular}
\end{table}

\subsection{Recuperación híbrida: BM25, similitud densa, RRF y el problema del medio}
Al recibir una consulta del usuario ---por ejemplo, ``¿Cuáles son los requisitos de elegibilidad para la ayuda al arranque de viñedo en el POSEI?''--- el sistema la proyecta al mismo espacio latente que los fragmentos indexados. El recuperador denso evalúa la similitud coseno entre el vector de la pregunta y los vectores de los chunks. Esa vía captura paráfrasis (``prima de arranque'' frente a ``ayuda al abandono de viña''), pero falla con identificadores alfanuméricos exactos (``Orden APA/102/2024'', códigos SIGPAC, números de anexo).

La vía léxica clásica, BM25 (Robertson y Zaragoza, 2009), puntúa coincidencia de términos con saturación de frecuencia y penalización por longitud del documento. En un corpus normativo es indispensable: el inspector busca a menudo un código, no un sinónimo. El prototipo mantiene un índice BM25 por tenant, reconstruido tras la ingesta, con el mismo universo de chunks que \texttt{pgvector}.

Fusionar dos rankings no es promediar puntuaciones: las escalas de coseno y de BM25 no son comparables. Reciprocal Rank Fusion (Cormack, Clarke y Buettcher, 2009) asigna a cada documento $1/(k+r)$ en cada lista y suma, con $k=60$ en este trabajo. El efecto es que un fragmento mediocre en ambas listas puede superar a uno excelente en una sola, y que un código capturado solo por BM25 sigue teniendo oportunidad de entrar al contexto. El Cross-Encoder (\texttt{bge-reranker-v2-m3}) relee el par pregunta--fragmento y reordena; en la demostración está desactivado por latencia, decisión empírica que se documenta en la configuración experimental.

Incluso con un top-$k$ correcto, el generador puede ignorar evidencia colocada en el centro del prompt (\emph{lost in the middle}; Liu et al., 2023). Por eso $k$ final no debe inflarse: más contexto no es más fidelidad si el modelo atiende sobre todo al principio y al final. El chat usa $k=8$; un $k=10$ de laboratorio es un compromiso, no una ley.

\subsection{Naive RAG, Graph RAG y RAG agéntico (literatura frente al prototipo)}
La literatura distingue, de forma esquemática, tres familias. El RAG ingenuo recupera y genera en un solo paso. El Graph RAG modela entidades y relaciones y recorre el grafo en la consulta. El RAG agéntico planifica, llama a herramientas y verifica. En AgroPS el chat híbrido recupera siempre; el modo agéntico orquesta un grafo LangGraph que decide si buscar o responder, puntúa la relevancia del contexto y puede reescribir la pregunta. No se evalúa un Graph RAG operativo: la proyección UMAP es visualización. Presentar esas familias aquí sitúa el prototipo; las métricas del Capítulo~\ref{ch:evaluacion} miden el retriever híbrido, no una ablación del bucle de reescritura.

Fase de generación condicionada: los $k$ fragmentos se ensamblan en una plantilla junto con la consulta y una instrucción de no invención. El modelo genera de forma autorregresiva, pero ahora la distribución de cada token está condicionada también al texto recuperado. Si ese texto es incompleto, el anclaje estricto produce abstención o respuesta parcial: mejor que una cifra inventada, insuficiente para un gestor que necesita el expediente completo.

\begin{table}[htbp]
\centering
\caption{Familias de RAG en la literatura y su relación con el prototipo.}
\label{tab:rag-tipos}
\begin{tabular}{p{3.2cm}p{5.4cm}p{5.4cm}}
\toprule
\textbf{Tipo} & \textbf{Idea} & \textbf{En este TFM} \\
\midrule
Naive RAG & Recuperar y generar en un paso & Base del chat híbrido (siempre recupera) \\
Modular RAG & Indexación, retrieve y generate como módulos sustituibles & \textbf{Arquitectura evaluada}: hexágono + denso + BM25 + RRF \\
Graph RAG & Entidades y aristas; multi-salto & No operativo; UMAP es visualización \\
Agentic RAG & Planificación y herramientas & En el producto (LangGraph); \textbf{no} es la corrida congelada \\
Multimodal RAG & Imagen/vídeo en la consulta & Solo OCR de ingesta; evaluación textual \\
\bottomrule
\end{tabular}
\end{table}

\textbf{Decisión de este TFM.} Se implementa un \textbf{RAG textual modular} (Gao et al., 2024): el recuperador y el generador son módulos con puertos sustituibles. El chat por defecto es naive-modular (un paso retrieve$\rightarrow$generate) con fusión híbrida. El modo agéntico se ofrece en la interfaz; las tablas del Capítulo~\ref{ch:evaluacion} miden el híbrido fijo. No se evalúa GraphRAG ni un RAG multimodal de consulta.

\section{Revisión crítica de trabajos previos}
En los últimos años han proliferado estudios sobre arquitecturas Retrieval-Augmented Generation (RAG) para mitigar alucinaciones en LLMs, especialmente en dominios regulados como la medicina y el derecho. Estas investigaciones coinciden en que RAG incorpora conocimiento externo ``no paramétrico'' (por ejemplo, bases de datos vectoriales o textos normativos) al proceso de generación, lo que reduce la tasa de respuestas erróneas o fabricadas. Por ejemplo, en entornos médicos especializados, Wada et al.\ (2025) \emph{reportaron} una tasa de alucinación nula en un conjunto acotado de consultas de radiología (0\,\% frente a 8\,\% del modelo base). Ese resultado no se transfiere a POSEI: dominio, $N$ y definición de alucinación son distintos, y este TFM \textbf{no} adopta la tesis de que RAG lleve las alucinaciones ``a niveles cercanos a cero''. De forma similar, He et al. (2025) mostraron que un modelo GPT-4 con RAG alcanzó un 96.4 \% de precisión en evaluaciones preoperatorias (vs. 86.6 \% humano) sin generar alucinaciones. En el ámbito jurídico, Schwarcz et al. (2025) hallaron que una herramienta jurídica basada en RAG mejoraba la calidad del trabajo (claridad y profesionalismo) y ``redujo las alucinaciones a niveles comparables'' al trabajo humano sin IA. Igualmente, Rabnam et al. (2025) desarrollaron el framework LQ-RAG para consultas legales, logrando mejoras significativas en métricas de recuperación (por ejemplo, +13 \% de hit rate y +15 \% de MRR) mediante agentes de auditoría y modelos afinados, aunque sin incorporar aún retroalimentación humana.

En agricultura, Sawant et al. (2026) diseñaron un sistema RAG para asesoría agronómica, procesando documentos de prácticas agrarias mediante chunking semántico y embeddings especializados. Evaluaron varios LLMs (Llama3.1, Mistral, Phi3, Qwen2.5) utilizando métricas de relevancia y recuperación (precisión@K, recall@K, MRR, NDCG) y métricas de calidad semántica (BLEU, ROUGE, BERTScore, faithfulness). Sus hallazgos indican que modelos como Mistral y Qwen2.5 obtienen mejor desempeño en relevancia y calidad de respuesta en el dominio agrícola. No obstante, casi ninguna investigación previa ha aplicado RAG explícitamente a la documentación normativa agraria (PAC/POSEI), lo que señala la brecha que ocupa este TFM.

\begin{table}[htbp]
\centering
\caption{Trabajos previos sobre RAG, alucinaciones, métricas y dominios regulados. Lectura crítica: qué se transfiere a este TFM y qué no.}
\label{tab:trabajos-previos}
\small
\begin{tabular}{p{2.3cm}p{2.0cm}p{2.4cm}p{3.4cm}p{3.6cm}}
\toprule
\textbf{Trabajo} & \textbf{Dominio} & \textbf{Arquitectura} & \textbf{Métricas} & \textbf{Crítica / hueco} \\
\midrule
Lewis et al.\ (2020) & QA general & RAG original & EM / F1 & Patrón fundacional; no normativa agraria. \\
Huang et al.\ (2025); Ji et al.\ (2023) & Transversal & --- & Taxonomía intrínseca/extrínseca & Define el constructo; no evalúa un retriever. \\
Es et al.\ (2024) & RAG genérico & --- & Faithfulness, answer relevancy, context precision/recall & Marco RAGAS; \textbf{no ejecutado} en la corrida congelada de este TFM. \\
Zhang et al.\ (2024) RAGTruth & Hallucinaciones en RAG & Varios RAG & Anotación humana & No contiene PAC/POSEI. \\
Wada et al.\ (2025) & Radiología & RAG local & Tasa de alucinación 0\,\% vs 8\,\% & $N$ y dominio no transferibles; no se adopta el ``cero alucinaciones''. \\
He et al.\ (2025) & Clínico & GPT-4 + RAG & Precisión vs humano & ``Sin alucinaciones'' es un resultado local, no una ley. \\
Schwarcz et al.\ (2025) & Jurídico & RAG legal & Calidad + hall.\ vs humano & ``Niveles comparables'' $\neq$ erradicación. \\
Rahman / LQ-RAG (2025) & Jurídico & Agéntico + fine-tune & Hit rate, MRR & Mejora IR; no es corpus POSEI. \\
Sawant et al.\ (2026) & Agronomía & RAG textual & P@$K$, MRR, faithfulness & Asesoría de prácticas, no ayudas PAC/POSEI. \\
Este TFM (2026) & PAC/POSEI Canarias & RAG textual modular híbrido & MRR, cobertura KW, juez 1--5, hall.\ binaria & Comparación pareada con/sin RAG; $N=35$; juez$\neq$humano. \\
\bottomrule
\end{tabular}
\end{table}

\section{Alucinaciones en modelos de lenguaje en entornos normativos (definición, causas y estudios relevantes).}
Los enfoques de Retrieval-Augmented Generation (RAG) combinan la recuperación de información externa con la generación de texto para enriquecer las respuestas de los LLM con evidencias verificables. La bibliografía resalta que RAG pretende mitigar las limitaciones paramétricas de los LLM, mejorando la calidad factual de las respuestas y reduciendo la incidencia de alucinaciones o afirmaciones no fundamentadas. En este contexto, una alucinación se entiende como una salida factualmente incorrecta que el modelo genera de manera plausible pero no está respaldada por la información recuperada. Sin embargo, varios autores advierten que la eficacia de RAG depende críticamente de la precisión de la etapa de recuperación: por ejemplo, Reuter et al. (2025) señalan que RAG sólo reduce las alucinaciones cuando la recuperación documental es fiable, advirtiendo fallos como el Document-Level Retrieval Mismatch donde se seleccionan textos irrelevantes. En suma, RAG no elimina automáticamente las alucinaciones; su beneficio debe verificarse empíricamente caso a caso.

La revisión de trabajos previos muestra diversas arquitecturas RAG especializadas en dominios regulados. Xu et al. (2025) introducen MEGA-RAG, un framework RAG para salud pública con múltiples fuentes de evidencia. Reportaron una reducción de alucinaciones de más del 40\% frente al modelo base y alta exactitud en las respuestas (precisión $\approx$0,79). En el ámbito legal, Rahman et al. (2025) proponen LQ-RAG, que integra agentes de auditoría y modelos finamente ajustados para consultas jurídicas. Su evaluación muestra mejoras relevantes en métricas de recuperación (por ejemplo, +13\% en Hit Rate, +15\% en MRR) y en la relevancia de las respuestas (hasta +23\% vs configuración ingenua). Feng et al. (2026) presentan Hyper-RAG, que emplea modelos de conocimiento basados en hipergrafos. En experimentos multiddominio (medicina, legal, agricultura) vieron un incremento promedio de precisión del 12.3\% sobre el LLM directo y mejoras del \textasciitilde{}6\% respecto a RAG con grafos tradicionales. Además, obtuvieron un +35.5\% de desempeño frente a RAG estándar en evaluación selectiva, confirmando la robustez de su método para consultas complejas. Por otro lado, Reuter et al. (2025) destacan que la fragmentación del texto legal puede obviar contexto global; proponen Summary-Augmented Chunking para enriquecer cada segmento con un resumen sintético, mejorando la recuperación (precisión/recall del contexto) y reduciendo las incoherencias de contexto. En el dominio agrícola normativo, Li et al. (2026) introducen AgriTune-R, un protocolo integral (adaptación LoRA, RAG, verificación experta) para LLM agrarios. Aunque es un marco metodológico sin resultados cuantitativos publicados, enfatiza la trazabilidad y consistencia de fuentes en explicaciones de políticas agrícolas.

Ninguno de estos trabajos afirma que RAG elimine las alucinaciones; todos tratan la reducción de información incorrecta como hipótesis a validar. Los hallazgos disponibles indican mejoras sustanciales en métricas de fidelidad o exactitud, pero dependen del dominio y de la calidad de los componentes de RAG. En muchos casos se carece de métricas estandarizadas y de pruebas con documentos normativos agrícolas concretos. Por ejemplo, el corpus RAGTruth (Zhang et al., 2024) provee un dataset para medir alucinaciones en RAG con anotaciones detalladas, pero no está especializado en la documentación PAC/POSEI. En conclusión, la literatura sugiere que RAG puede aumentar la fidelidad de las respuestas, pero subraya la necesidad de evaluarlo rigurosamente en el contexto normativo agrícola mediante métricas cuantitativas.

\section{Métricas de evaluación: definiciones y lo que mide este prototipo}
La evaluación de un RAG exige separar recuperación y generación (Es et al., 2024). Las definiciones siguientes son las de la literatura. \textbf{En la corrida congelada} se calculan de hecho: MRR y cobertura de palabras clave (funciones propias); exactitud, exhaustividad y relevancia 1--5 más \texttt{is\_hallucination} (juez local). Faithfulness, answer relevancy, context precision y context recall de RAGAS \textbf{no} se reportan como cifras oficiales porque la librería no se ejecutó. La calidad de citas se discute cualitativamente (metadatos de fragmento en el chat) y con fallos de identificador en el Capítulo~\ref{ch:evaluacion}.

\paragraph{Fidelidad (faithfulness).}
Sea $A$ la respuesta y $\{a_i\}_{i=1}^{n}$ las afirmaciones extraídas de $A$. Si $n_f$ de ellas se infieren del contexto recuperado $C$,
\begin{equation}
\mathrm{Faithfulness}(A,C)=\frac{n_f}{n}\in[0,1].
\end{equation}
En este TFM el proxy empírico es la marca binaria \texttt{is\_hallucination} del juez (respuesta no sustentada) y, de forma inversa, la exactitud 1--5 respecto a la referencia.

\paragraph{Relevancia de la respuesta (answer relevancy).}
Se estima la similitud entre la pregunta $q$ y la respuesta $A$ (en RAGAS, generando preguntas inversas y promediando el coseno de embeddings). Aquí el proxy es la nota \texttt{relevance} del juez (1--5).

\paragraph{Precisión del contexto (context precision).}
Si $\mathrm{rel}(k)\in\{0,1\}$ indica si el fragmento en posición $k$ es relevante,
\begin{equation}
\mathrm{ContextPrecision@}K=\frac{1}{|\mathcal{R}|}\sum_{k=1}^{K}\mathrm{Precision@}k\cdot\mathrm{rel}(k),
\end{equation}
con $\mathcal{R}$ el conjunto de relevantes en el top-$K$. En la corrida se aproxima con MRR y cobertura de palabras clave sobre el ranking híbrido.

\paragraph{Exhaustividad del contexto (context recall).}
Si $G$ es el conjunto de afirmaciones de la respuesta de referencia y $G_C\subseteq G$ las cubiertas por $C$,
\begin{equation}
\mathrm{ContextRecall}=\frac{|G_C|}{|G|}\in[0,1].
\end{equation}
El proxy en este trabajo es la cobertura de keywords del ítem en los chunks recuperados ($0{,}9643$ en la corrida).

\paragraph{Exactitud factual.}
Acuerdo de la respuesta generada con la referencia de oro. Proxy: \texttt{accuracy} del juez (1--5), pareada con/sin RAG (Tabla~\ref{tab:rag-vs-norag}).

\paragraph{Tasa de alucinaciones.}
\begin{equation}
\mathrm{HallucinationRate}=\frac{1}{N}\sum_{j=1}^{N}\mathbf{1}[\text{la respuesta }j\text{ no está sustentada}].
\end{equation}
Operacionalización: fracción de ítems con \texttt{is\_hallucination}=\texttt{true} según el juez.

\paragraph{Calidad de las citas.}
Precisión de cita: proporción de referencias (documento, sección/página, fragmento) que coinciden con evidencia real. Cobertura de cita: fracción de la referencia de oro apoyada por al menos una cita. En el prototipo cada fuente arrastra \texttt{source\_file}, sección (\texttt{page} del banco) e identificador de chunk; la página PDF exacta depende del parser. Los fallos H1 y H2 muestran que la UI puede citar un documento vecino mientras el texto inventa un código: la citación visible \textbf{no} equivale a citación correcta.




